\section{Low-depth OPE constraints for $V(\mathfrak{g}_{\Delta_5})$}
\label{wn:sec:delta5-vacuum-OPE-depth-5}

\providecommand{\gDelta}{\mathfrak{g}_{\Delta_5}}
\providecommand{\VDelta}{V(\gDelta)}
\providecommand{\VLatt}{V_{\Lambda^{2,1}}}
\providecommand{\NO}[1]{\mathopen{:}{#1}\mathclose{:}}
\providecommand{\kBKM}{\kappa_{\mathrm{BKM}}}
\providecommand{\kch}{\kappa_{\mathrm{ch}}}
\providecommand{\kcat}{\kappa_{\mathrm{cat}}}
\providecommand{\kfib}{\kappa_{\mathrm{fiber}}}

\paragraph{Scope.}
Assume a conformal vertex superalgebra $\VDelta$ whose Borcherds
denominator is the Gritsenko--Nikulin product $\Delta_5$ and whose
Virasoro field has central charge $c=-214$. The OPEs below are
constraints on any screened free-field model for $\VDelta$. They do not
construct the model.

\paragraph{Virasoro field.}
The class-$\mathcal S$ parent
$\mathcal T[A_1,\Sigma_{0,24}]$ has
$(n_v,n_h)=(63,88)$, hence
\[
  c_{4d}={2n_v+n_h\over 12}={107\over 6},
  \qquad
  c_{2d}=-12c_{4d}=-214
\]
by the Beem--Rastelli four-dimensional to two-dimensional anomaly
formula. If $T(z)$ is the stress tensor of the corresponding protected
vertex algebra, then the Virasoro OPE is
\begin{equation}
\label{wn:eq:delta5-TT-ope}
T(z)T(w)
 =
 { -107\over (z-w)^4}
 + {2T(w)\over (z-w)^2}
 + {\partial T(w)\over z-w}
 + \mathrm{reg}.
\end{equation}
This fixes only the Virasoro anomaly. It does not identify the full
operator algebra with a Borcherds denominator module.

\paragraph{Cartan currents and real-root fields.}
Let $\delta_1,\delta_2,\delta_3$ be the real simple roots of the
rank-three Lorentzian chamber of $\gDelta$, with Gram matrix
\[
  (\delta_i,\delta_j)=
  \begin{cases}
  2, & i=j,\\
  -2, & i\ne j.
  \end{cases}
\]
For free bosons $\varphi$ valued in this lattice, the Heisenberg
currents $H_i(z)=\NO{\delta_i\cdot\partial\varphi}(z)$ satisfy
\begin{equation}
\label{wn:eq:delta5-Heisenberg-ope}
H_i(z)H_j(w)
 =
 {(\delta_i,\delta_j)\over (z-w)^2}
 + \mathrm{reg}.
\end{equation}
The lattice exponentials
$E_i(z)=\NO{\exp(\delta_i\cdot\varphi)}(z)$ obey the lattice OPE
\begin{equation}
\label{wn:eq:delta5-lattice-exp-ope}
E_i(z)E_j(w)
 =
 \varepsilon_{ij}(z-w)^{(\delta_i,\delta_j)}
 \NO{\exp((\delta_i+\delta_j)\cdot\varphi)}(w)
 + \cdots .
\end{equation}
Thus the affine-current formula applies to the Cartan currents
$H_i$, while the real-root fields carry the usual lattice singularities.
For $i\ne j$, the product has a double pole and lands in the null
charge $\delta_i+\delta_j$. A Borcherds vertex realization must impose
the correct quotient or screening differential on these null-charge
fields.

\paragraph{Imaginary-root multiplicity.}
In the Gritsenko--Nikulin normalization used for $\Delta_5$, the Jacobi
input has constant coefficient $c(0)=10$ and polar coefficient
$c(-1)=1$. The first even null multiplicity is therefore
\[
  \tau(a_0)=c(0)-c(-1)=9 .
\]
This integer is a Borcherds root multiplicity. A conformal weight for a
field of charge $a_0$ depends on the chosen free-field background charge
and ghost realization. The number $9$ supplies no such weight by itself.

\paragraph{Virasoro degeneracy at depth five.}
For the Virasoro Verma determinant write
\[
  c=13-6(t+t^{-1}),\qquad
  h_{r,s}(c)=
  { (rt-s)^2-(t-1)^2 \over 4t } .
\]
At $c=-214$, the parameter $t$ satisfies
$6t^2-227t+6=0$, so $t$ is irrational. The equation
$h_{r,s}(-214)=0$ with $rs\le 5$ forces $(r,s)=(1,1)$.
After quotienting the translation vector $L_{-1}|0\rangle$, there is
no further Virasoro singular vector through depth five. Any low-depth
null relation in $\VDelta$ must therefore come from the Borcherds
screening data, not from a Virasoro minimal-model degeneration.

\paragraph{Screening model.}
A Feigin--Frenkel style presentation
\[
  \VDelta \simeq H^\bullet(\VLatt\otimes \mathcal G,Q)
\]
requires a ghost system $\mathcal G$ and a BRST operator $Q$ whose
screening currents all have conformal weight one, mutually local OPEs,
and square-zero total charge. The proof obligations are explicit:
construct the real-root screenings, construct the null and odd
imaginary screenings with multiplicities read from the Jacobi
coefficients, prove $Q^2=0$, compute the cohomology, and match the
graded supercharacter with the Borcherds product for $\Delta_5$.
Central charge $-214$ and the scalar $\tau(a_0)=9$ do not imply this
screening theorem.

\paragraph{Invariant separation.}
For $K3\times E$ the compact Hodge supertrace is
$\kch(K3\times E)=0$, while the Heisenberg specialization has
$\kch^{\mathrm{Heis}}(K3\times E)=3$. The categorical invariant is
$\kcat(K3\times E)=0$. The Borcherds invariant is
$\kBKM(\Delta_5)=c_1(0)/2=10/2=5$. The fibre invariant is
$\kfib(K3)=24$. These are four distinct constructions.
